\color{unidarkgreen}{\subsection{Логика опробования}}
\color{black}

\begin{enumerate}

\item
Логика опробования предназначена для работы ДЗО при оперативном опробовании выключателей присоединений по схеме с «открытым» или «закрытым» плечом. При неуспешном опробовании отключается только опробуемый выключатель, остальные выключатели остаются в работе.

\item
Алгоритм работы ЛОП приведен на рисунке \ref{pic_log_oprob:Log_oprob}. Уставки ЛОП приведены в таблице \ref{app_set:lop}.

\medskip
\begin{figure}[H]
\centering
\includegraphics[width=0.83\textwidth,height=0.83\textheight,keepaspectratio]{\fbpath/05.01 ДЗО/661.1341 Логика опробования/img/Log_oprob.pdf}
\captionof{figure}{Функциональная схема ЛОП}\label{pic_log_oprob:Log_oprob}
\end{figure}
\medskip

\item
Переход в режим опробования осуществляется при наличии разрешающего сигнала \textit{«Ввод опробования»}. Запуск опробования выполняется при поступлении команды РКВ от выключателя любого присоединения. Продолжительность опробования определяется уставкой выдержки времени \textbf{«Топроб»}.

\item
С помощью программного переключателя \textbf{«Опроб. Вn»} задается условие для отключения опробуемого выключателя по срабатыванию одного из следующих органов: ДЗО, ЧТО или индивидуального ИО по току опробуемого присоединения с блокировкой по второй гармонике.

\item
Ток срабатывания индивидуального ИО задается с помощью уставки \textbf{«Iср опроб. Вn»}. Блокировка по второй гармонике необходима для предотвращения ложных срабатываний ИО при БНТ трансформатора. Для ввода блокировки предусмотрен программный ключ \textbf{«БНТ Вn»}. Уровень срабатывания блокировки задается с помощью уставки \textbf{«I2г/I1г Вn»}. Срабатывание ИО осуществляется по завершении выдержки времени \textbf{«Тср. опроб. Bn»}.

\item
Для обеспечения отключения только опробуемого выключателя, без отключения всей СШ, предусмотрены сигналы блокировки ДЗО: \textit{«ЛОП: Блок. ДЗО при опроб.»}. Время действия блокировки соответствует продолжительности опробования.


\end{enumerate}

